
\documentclass{article}
\usepackage{colortbl}
\usepackage{makecell}
\usepackage{multirow}
\usepackage{supertabular}

\begin{document}

\newcounter{utterance}

\centering \large Interaction Transcript for game `bbh', experiment `boolean\_expressions', episode 71 with qwen3.5{-}4b.
\vspace{24pt}

{ \footnotesize  \setcounter{utterance}{1}
\setlength{\tabcolsep}{0pt}
\begin{supertabular}{c@{$\;$}|p{.15\linewidth}@{}p{.15\linewidth}p{.15\linewidth}p{.15\linewidth}p{.15\linewidth}p{.15\linewidth}}
    \# & \multicolumn{2}{c}{Player} && \multicolumn{2}{c}{Game Master} \\
    \hline

    \theutterance \stepcounter{utterance}  
    & & & \multicolumn{4}{p{0.6\linewidth}}{
        \cellcolor[rgb]{0.9,0.9,0.9}{
            \makecell[{{p{\linewidth}}}]{
                \texttt{\tiny{[P1$\langle$GM]}}
                \texttt{Evaluate the result of a random Boolean expression.} \\
\\ 
\texttt{Q: not ( ( not not True ) ) is} \\
\texttt{A: False} \\
\\ 
\texttt{Q: True and False and not True and True is} \\
\texttt{A: False} \\
\\ 
\texttt{Q: not not ( not ( False ) ) is} \\
\texttt{A: True} \\
\\ 
\texttt{Q: not not not not not not not True is} \\
\texttt{A:} \\
            }
        }
    }
     \\ \\

    \theutterance \stepcounter{utterance}  
    & \multicolumn{4}{p{0.6\linewidth}}{
        \cellcolor[rgb]{0.561,0.737,0.561}{
            \makecell[{{p{\linewidth}}}]{
                \texttt{\tiny{[P1$\rangle$GM]}}
                \texttt{To evaluate the expression `not not not not not not not True`, let's break it down step{-}by{-}step.} \\
\\ 
\texttt{The expression consists of 7 `not` operators applied to the value `True`.} \\
\texttt{Recall that:} \\
\texttt{{-} `not True` = `False`} \\
\texttt{{-} `not False` = `True`} \\
\texttt{{-} Applying `not` twice returns the original value (`not not X` = `X`).} \\
\\ 
\texttt{We can group the operators in pairs:} \\
\texttt{1.  **First pair**: `not not` cancels out.} \\
\texttt{2.  **Second pair**: `not not` cancels out.} \\
\texttt{3.  **Third pair**: `not not` cancels out.} \\
\texttt{4.  **Remaining operator**: There is one `not` left.} \\
\\ 
\texttt{Let's count the total number of `not` operators:} \\
\texttt{{-} Total `not`s = 7} \\
\texttt{{-} Pairs of `not`s = 3 pairs (6 operators)} \\
\texttt{{-} Remaining `not`s = 1} \\
\\ 
\texttt{So, the expression simplifies to:} \\
\texttt{`not ( (not not (not not (not not True))) )`} \\
\\ 
\texttt{Since `not not True` = `True`, we can reduce the inner parts:} \\
\texttt{`not ( (not not (not not True)) )`} \\
\texttt{`= not ( (not not True) )`} \\
\texttt{`= not ( True )`} \\
\texttt{`= False`} \\
\\ 
\texttt{Alternatively, simply counting the parity of the `not` operators:} \\
\texttt{{-} An odd number of `not`s flips the value.} \\
\texttt{{-} An even number of `not`s keeps the value.} \\
\texttt{{-} 7 is an odd number.} \\
\texttt{{-} Input: `True`} \\
\texttt{{-} Output: `False`} \\
\\ 
\texttt{A: False} \\
            }
        }
    }
    & & \\ \\

    \theutterance \stepcounter{utterance}  
    & & & \multicolumn{2}{p{0.3\linewidth}}{
        \cellcolor[rgb]{0.9,0.9,0.9}{
            \makecell[{{p{\linewidth}}}]{
                \texttt{\tiny{[GM$|$GM]}}
                \texttt{Target: False} \\
            }
        }
    }
    & & \\ \\

    \theutterance \stepcounter{utterance}  
    & & & \multicolumn{2}{p{0.3\linewidth}}{
        \cellcolor[rgb]{0.9,0.9,0.9}{
            \makecell[{{p{\linewidth}}}]{
                \texttt{\tiny{[GM$|$GM]}}
                \texttt{game\_result = LOSE} \\
            }
        }
    }
    & & \\ \\

\end{supertabular}
}

\end{document}
